\section{Rigidity under a finite density bound}\label{sec:finite}

We now prove Theorem~\ref{thm:finite}, in its full range $4\le n\le7$.
This will also make a density bound on one sufficiently large ball a
regularity criterion; that local consequence is the input for the later
analysis without a global density bound. Write $n=m+1$, let $\omega_m$
be the volume of the unit ball in $\R^m$, and keep
$s_0=\sigma/2=2\sqrt2/3$. Recall the normalized energy

\[
 M_r(y)=\frac{1}{s_0\omega_m r^m}\int_{B_r(y)}e(v),\qquad
 e(v)=\frac12|\nabla v|^2+W(v),\qquad
 M_\infty=\sup_{r>0}M_r(0),
\]
as defined in \eqref{eq:finite-density}.

The statements from the literature used below are preliminary results of
Florit-Simon--Serra, with their numbering fixed to
\cite{FSS2025}, version~1.  Their final classification theorem concerns
\(n=4\); the extension to \(n=5,6,7\) here is the gap argument following
Lemma~\ref{lem:finite-critical}.

Suppose the theorem fails. The literature reduction produces a critical
solution with an integer number $K\ge2$ of zero layers outside a sparse
set. We will compare two estimates for the nearest gap between those
layers. A Green test bounds its average in terms of mean curvature;
Lemma~\ref{lem:diffuse-mean} and a Hardy test bound that curvature cost.
The opposite orientations of adjacent transitions force almost all
gaps to be logarithmically large. A preliminary square-root gap test is
needed before the sharper, nearly linear test: it makes the gradient
error in the gap Laplacian absorbable without losing the final coefficient.

\subsection{Reduction to a critical integer-density solution}

By \eqref{eq:finite-monotonicity}, $M_r(y)$ is nondecreasing and
its limit at infinity is independent of $y$. In the diffuse variables,
the discrepancy in that formula is $q^2D/4$ by
\eqref{eq:diffuse-energy}; thus
\[
 M_r'(y)=\frac{1}{s_0\omega_m}
 \left\{r^{-m-1}\int_{B_r(y)}\frac{q^2D}{4}
       +r^{-m}\int_{\partial B_r(y)}|\partial_rv|^2\right\}.
\]
We will use its two nonnegative terms separately. The derivation in
Section~\ref{sec:classical-estimates} uses neither stability nor finite
density; finite density now enters through the following reduction.

A number \(K_f>0\) is called \emph{subcritical} in dimension \(n\) if the
conclusion of Theorem~\ref{thm:finite} holds for every bounded entire stable
solution with \(M_\infty\le K_f\).  The range of a bounded entire solution
is automatically contained in \([-1,1]\), by applying the equation to
translated limits at its supremum and infimum.  The constants \(\pm1\)
are included in this definition.

\paragraph{The two regularity inputs.}
There are two different ways to obtain a graphical region: subcritical
energy density, or small local generalized-curvature mass. We state
both hypotheses because the critical solution uses the second before
the theorem makes the first available at every finite threshold. For a regular point put
\(\mathcal A_v=|\nabla(\nabla v/|\nabla v|)|\), where the derivative is the
full ambient derivative.  The \emph{sheeting assumptions} on \(B_r\) are
\[
 |\nabla v|\ge C^{-1},\qquad \mathcal A_v\le C/r
 \quad\hbox{on }B_r\cap\{|v|\le0.9\}.
\]
Thus a tangential curvature bound alone would not be this hypothesis.
For entire stable parameter-one solutions in \(n\le7\),
\cite[Theorem~3.2]{FSS2025} gives these assumptions on \(B_{\delta r}\)
when \(K_f\) is subcritical, \(M_r\le K_f+\delta\), and \(r^{-1}\le\delta\);
\(\delta>0,C<\infty\) depend on \(K_f,n\).
The separate good-ball result \cite[Theorem~3.4]{FSS2025} assumes
\(M_r\le C_0\), \(r\ge r_0(C_0,n)\), and
\[
 \int_{B_1(x)}\mathcal A_v^2|\nabla v|^2<\delta_{\rm bad}
 \quad\hbox{for every }x\in B_r\cap\{|v|\le0.9\}.
\]
It gives the sheeting assumptions on \(B_{r/2}\), with constants depending
on \(C_0,n\).  In this integral the curvature density is defined almost
everywhere as
\(
 \mathcal A_v^2|\nabla v|^2
 =|\nabla^2v|^2-|\nabla|\nabla v||^2.
\)

\paragraph{What sheeting supplies.}
Under the sheeting assumptions, \cite[Lemma~2.13 and Theorem~2.14]{FSS2025}
give graphical decomposition and the Wang--Wei estimates
\cite{WangWei2019}; the latter permit \(n\le10\).
Theorem~1.1 of the final Wang--Wei author manuscript, dated
9 October 2019, uses a potential with second derivative one at the
wells.  Our normalization is transferred to theirs by
\[
 \widehat W=W/2,\qquad \widehat v(y)=v(y/\sqrt2),\qquad
 \Delta_y\widehat v=\widehat W'(\widehat v).
\]
Then \(\widehat W''(\pm1)=1\), stability is preserved by change of
variables, and lengths are multiplied by \(\sqrt2\).
This fixed rescaling changes only the constants in the estimates below;
the parameter-one formulation in \cite[Theorem~2.14]{FSS2025}
already uses our quartic normalization.
On smaller cylinders they give scaled \(C^{2,\theta}\) bounds, mean
curvature \(O(r^{-2})\), and separation tending to infinity as
\(r\to\infty\).  We will only need the weaker
\(C_\iota r^{-2+\iota}\) mean-curvature bound for each fixed \(\iota>0\).
When all zeros in the original ball lie in a sufficiently thin slab,
\cite[Theorem~4.8 and Lemmas~4.9--4.10]{FSS2025} give arbitrarily small
slopes, comparison of density with the integer number of graphs, and
local approximation by a planar heteroclinic on every fixed window.
Interior elliptic estimates upgrade that approximation to smooth local
convergence.  The thickness threshold, the interior fractions, and the
constants are uniform once the sheeting constants, dimension, and desired
accuracy have been fixed.

\Needspace{3\baselineskip}
\begin{lemma}[Critical solution and its good regions]\label{lem:finite-critical}
If Theorem~\ref{thm:finite} fails in a fixed dimension \(4\le n\le7\),
there is a nonplanar entire stable solution with finite integer density
\(M_\infty=K\ge2\), and a nonempty closed set \(Z\), with the following
properties.  Put \(s(x)=\dist(x,Z)\) and \(\Sigma=\{v=0\}\).
\begin{enumerate}[label=\textup{(\roman*)}]
\item Uniformly for \(y\in Z\),
\[
 M_r(y)\longrightarrow K,\qquad
 r^{2-m}|B_1(Z)\cap B_r(y)|\longrightarrow0
 \quad(r\to\infty).
\]
The second conclusion holds with \(B_h(Z)\) in place of \(B_1(Z)\)
for each fixed \(h<\infty\).
\item There are \(c>0,C<\infty\) and uniform moduli tending to zero
as \(s\to\infty\) such that, about each \(z\in\Sigma\) with \(s(z)\)
large, the zero set in an interior cylinder of radius \(c s(z)\)
consists of exactly \(K\) ordered smooth graphs.  They are defined on
a disk of that radius over a hyperplane through \(z\), their heights
are \(o(s(z))\), their slopes are \(o(1)\), and
\[
 \|\nabla^2 f_i\|_\infty+
 s(z)^{1/2}[\nabla^2 f_i]_{C^{1/2}}\le C/s(z),
 \qquad |H_i|\le C_\iota s(z)^{-2+\iota}.
\]
Their separations tend uniformly to infinity.  At the inner zero
points \(D=o(1)\), and every fixed normal window converges smoothly
to the oriented \(g\)-profile.
\item Uniformly for \(y\in Z\),
\begin{equation}\label{eq:finite-discrepancy}
 r^{-m}\int_{B_r(y)}q^2D\longrightarrow0.
\end{equation}
\end{enumerate}
Constants and moduli in this lemma may depend on the fixed critical
solution and dimension.
\end{lemma}

\begin{proof}
We first obtain the critical solution and uniform information at its bad
centers. We then transfer that information to every distant good zero,
where the later gap test will be applied.

\emph{Attainment and the bad set.}
Near-unit-density rigidity, stated in \cite[Theorem~2.16]{FSS2025}
and originating in \cite[Theorem~11.1]{Wang2017}, makes
\(1+\delta(n)\) subcritical.  If some finite density is not subcritical,
the finite critical density is attained by
\cite[Proposition~3.6]{FSS2025}.  Take the bad set in its
Definition~3.7: its centers satisfy \(|v|\le0.9\) and the reverse
nonstrict inequality in the good-ball curvature condition above.
Lemma~3.8 there makes this set nonempty.
It is harmless to take its closure; all estimates below pass to that
closure by continuity and shifted-ball containment.
Propositions~3.9--3.10 there give uniform large-scale planar confinement,
uniform convergence of density at these centers, integrality of the
critical density, and the no-gaps conclusion.  Thus \(K\ge2\).
Lemma~5.1 and Proposition~5.2 there give the asserted sparse bad-set
estimate.  These results explicitly assume \(n\le7\), including
Section~5 of that paper.

For clarity, the enlargement from \(1\) to a fixed \(h\) costs only a
constant depending on \(h,n\).  Choose a maximal separated family of
bad centers in \(B_{r+h+2}(y)\) with fixed small separation.
Disjoint balls of a still smaller fixed radius about its centers lie
in \(B_1(Z)\cap B_{r+h+3}(y)\).  Balls of radius \(h+2\) about this
family cover \(B_h(Z)\cap B_r(y)\).  This proves the claim using
the estimate at the slightly larger radius.

\emph{Exactly $K$ sheets at every distant good zero.}
We now prove (ii), with the same moduli at all good points.  Choose a nearest
\(y\in Z\) to \(z\), with \(s=|z-y|\).  There are no bad centers in
\(B_{s/2}(z)\), so the good-ball theorem gives the sheeting assumptions
on \(B_{s/4}(z)\).  The plane supplied by large-scale flatness at
\(y\), at radius \(4s\), passes within \(o(s)\) of \(z\), and the zero
set in this ball is confined to an \(o(s)\)-slab about it.  Translate
the plane to pass through \(z\).

Here is the precise no-gaps transfer.  Write \(A_0=4s\), let
\(P_{y,A_0}\) be the plane through \(y\) given by large-scale
flatness, and let \(a_0\) be the orthogonal projection of \(z\)
to that plane.  There is one modulus \(\omega\), independent of
the bad center \(y\), such that
\[
 |z-a_0|\le A_0\omega(A_0^{-1}),\qquad
 M_{A_0\omega(A_0^{-1})}(a_0)\ge K-\omega(A_0^{-1}).
\]
The second inequality is exactly the conclusion of
\cite[Proposition~3.10]{FSS2025}; it applies because
\(a_0\in P_{y,A_0}\cap B_{A_0}(y)\).
For any fixed \(\theta>0\), the number
\(b_s=\theta s-|z-a_0|\) exceeds
\(A_0\omega(A_0^{-1})\) for all sufficiently large \(s\).
Monotonicity at \(a_0\) and \(B_{b_s}(a_0)\subset B_{\theta s}(z)\)
then give
\[
 K\ge M_{\theta s}(z)
 \ge\left(\frac{b_s}{\theta s}\right)^m
       \bigl(K-\omega(A_0^{-1})\bigr)=K-o(1).
\]
All thresholds and errors are uniform in \(z,y\); no diagonal
selection of centers or radii is involved.

Set \(a=s/4\), the radius of the ball where the sheeting assumptions
hold.  Translating \(P_{y,A_0}\) to pass through \(z\) confines every
zero in \(B_a(z)\) to a slab of width at most
\(2A_0\omega(A_0^{-1})=o(a)\).
Apply \cite[Theorem~4.8]{FSS2025} on \(B_a(z)\).
It represents all zeros in the cylinder
\(B'_{3a/5}\times[-3a/5,3a/5]\), in the resulting coordinates,
by \(N\) ordered graphs defined over \(B'_{3a/5}\), each with
height in \([-a/8,a/8]\).  Thus every listed graph meets
\(B_{a/2}(z)\), and no zero in that ball is omitted.
Its Lemma~4.9 gives, for any prescribed sufficiently small
\(\delta_1>0\) and all sufficiently large \(s\),
\[
 |M_{a/2}(z)-N|\le\delta_1N.
\]
The preceding density transfer with \(\theta=1/8\) gives
\(M_{a/2}(z)\to K\).  First \(\delta_1\le1/2\) gives
\(N\le2K\); then choose \(\delta_1<1/(8K)\) and a density
deficit smaller than \(1/4\).  It follows that
\(|N-K|<1/2\), and hence \(N=K\).
The curvature, separation and profile conclusions are those of the
same results.  All applications use fixed interior fractions and
uniform moduli.

\emph{Uniform disappearance of discrepancy.}
Integrate the first term of \eqref{eq:finite-monotonicity}
between \(r\) and \(2r\).  Its nonnegativity gives
\[
 r^{-m}\int_{B_r(y)}q^2D
 \le C\bigl(M_{2r}(y)-M_r(y)\bigr).
\]
Uniform convergence at bad centers proves (iii).
\end{proof}

For the rest of the proof fix such a critical solution.  The preliminary
rigidity case \(D=0\) is excluded, so \(D>0\).
For a sufficiently large fixed \(S\), write
\(\Sigma_S=\Sigma\cap\{s>S\}\).  This is a smooth open surface.
Every small quantity denoted below by \(\varepsilon_S\) is uniform and
tends to zero as \(S\to\infty\).  No constant depending on a radial
exhaustion scale or a cap will be hidden in this notation.

\subsection{Surface measures and radial estimates}

We first transfer the ambient energy and stability bounds to the regular
zero surface.  Four estimates will be used: defect mass grows at most
like \(r^{m-2}\), radial geometric errors have a vanishing tail, surface
mass retains the coefficient \(K\), and a Hardy test fixes the sharp
coefficient of the logarithmic defect integral.  All radial tests are
centered at arbitrary points of \(Z\).

\begin{lemma}[Tube, mass and Hardy estimates]\label{lem:finite-measures}
For any prescribed finite \(P\), increasing \(S\) makes
\[
 \Phi(z,t)=z+t\nu(z),\qquad z\in\Sigma_{S/2},\quad |t|\le P,
\]
injective, with Jacobian \(1+\varepsilon_S\) and
\(q(\Phi(z,t))=G(t)+o_S(1)\), uniformly.
Good-surface area in every ball of radius \(r\ge1\) is at most \(Cr^m\),
and
\begin{equation}\label{eq:B1}
 \int_{\Sigma_{S/2}\cap B_r(x)}D\,\dd S
 \le C\int_{B_{r+1}(x)}q^3D\,\dd x
 \le Cr^{m-2}\qquad(x\in\R^n,\ r\ge1).
\end{equation}
For \(y\in Z\), put \(r=|z-y|\) on the surface and
\(e_r=(z-y)/r\).  Uniformly in \(y\),
\begin{equation}\label{eq:B2}
 \int_{\Sigma_S\cap\{r>R_0\}}
 \left((e_r\cdot\nu)^2r^{-m}+|H|r^{1-m}\right)\dd S
 \le \eta(R_0),\qquad \eta(R_0)\longrightarrow0.
\end{equation}
The same function \(\eta\) may be used for every \(S\) above one
fixed unit-tube threshold \(S_*\).
Choose a smooth ambient \(0\le\chi\le1\), supported in \(\{s>S\}\)
and equal to one on \(\{s\ge2S\}\), with bounded first two derivatives.
For fixed nonnegative Lipschitz \(F\), supported in \([1/2,2]\), and
every \(\delta>0\), \(S\) can be chosen so that
\begin{equation}\label{eq:B3}
 \liminf_{R\to\infty}R^{-m}
 \int_{\Sigma}\chi F(r/R)\,\dd S
 \ge K m\omega_m\int_0^\infty t^{m-1}F(t)\,\dd t-\delta.
\end{equation}
Finally, for each fixed \(R_0\ge10\),
\begin{equation}\label{eq:B4}
 \limsup_{R\to\infty}\frac{1}{\sigma\log R}
 \int_{\{R_0/2<|x-y|<4R\}}q^3D\,|x-y|^{2-m}\,\dd x
 \le K m\omega_m\frac{(m-2)^2}{4}.
\end{equation}
The limiting estimates are uniform for \(y\in Z\).
\end{lemma}

\begin{proof}
We prove the assertions in their order: first tube injectivity and defect
mass, then the radial tail, then the surface mass retained outside the
bad set, and finally the Hardy coefficient.

\emph{Tube injectivity and defect mass.}
Curvature \(O(1/s)\) and profile convergence give the local tube
description.  To verify injectivity, a collision would have base
points within distance \(2P\), hence in a common good chart when
\(S\) is large.  Different sheets cannot collide because their separation
tends to infinity; on a single sheet the small slope and curvature
give uniqueness of the normal projection on this fixed-width tube.
The volume Jacobian is
\(\det(\Id+tA)=1+o_S(1)\).
A unit tube, on which \(q\) is bounded below, controls surface area by
energy.  Bounded density therefore gives the area bound for arbitrary
centers.  Local Harnack comparison in \eqref{eq:diffuse-local} also gives
the first inequality of \eqref{eq:B1}.  To obtain its second inequality,
apply \eqref{eq:q-stability} with a cutoff equal to one on \(B_{r+1}(x)\),
supported in \(B_{2r+2}(x)\), and with gradient bounded by \(C/r\);
use \(q^2\le4e(v)\).  On good zeros the identity
\(H=\partial_\nu D/(2|p|^2)\) gives \(|H|\le CD\).

\emph{The radial error has a vanishing tail.}
Integrating the second term in \eqref{eq:finite-monotonicity} gives
\[
 \int_{\{|x-y|>R_0\}}|\partial_rv|^2|x-y|^{-m}\,\dd x
 \le s_0\omega_m\bigl(K-M_{R_0}(y)\bigr).
\]
In a unit tube, \(|p(\Phi(z,t))-\nu(z)|\le C\sqrt{D(z)}\),
by \eqref{eq:diffuse-local}; changing radial directions costs \(C/r\).
After squaring, integrating in \(t\), and using the unit tube's positive
lower profile bound, the surface radial-normal integral is bounded
by this ambient flux tail, plus
\[
 C\int_{\Sigma_S\cap\{r>R_0\}}
       \bigl(Dr^{-m}+r^{-m-2}\bigr)\,\dd S.
\]
On a dyadic shell of radius \(a\), \eqref{eq:B1} bounds the first
error by \(Ca^{-2}\), and the area bound does the same for the second.
The mean-curvature term in \eqref{eq:B2} is at most \(Ca^{-1}\)
on that shell.  Summing proves the explicit uniform bound
\[
 \eta(R_0)\le C\left(
   \sup_{y\in Z}\bigl(K-M_{R_0-2}(y)\bigr)+R_0^{-1}\right),
\]
which tends to zero by uniform density convergence at \(Z\).
The constants can be fixed when the unit tube and
\eqref{eq:B1} are first valid, at \(S=S_*\).
Increasing \(S\) merely restricts the nonnegative integrals to
smaller subsets, so it does not worsen this bound.

\emph{Removing the bad set while retaining the coefficient $K$.}
We first construct the stated smooth cutoff without differentiating the
nonsmooth distance $s$.  Mollify \(s\) at radius
\(S/8\), writing \(\widetilde s=\rho_{S/8}*s\), and choose a smooth
nondecreasing \(\theta:\R\to[0,1]\) equal to zero on
\(( -\infty,5/4]\) and to one on \([7/4,\infty)\).
Then \(\chi=\theta(\widetilde s/S)\) has the required properties:
\(|\widetilde s-s|\le S/8\), so its derivatives are supported in
\(9S/8\le s\le15S/8\), strictly inside the regular region \(s>S\).
Moreover \(|\nabla\chi|\le C/S\) and
\(|\nabla^2\chi|\le C/S^2\).

The same tube and the fixed-thickness bad-set estimate imply
\begin{equation}\label{eq:finite-collar-area}
 \sup_{y\in Z}\cH^m\bigl(\Sigma\cap\{S<s<2S\}\cap B_r(y)\bigr)
   =o(r^{m-2})
 \quad(r\to\infty)
\end{equation}
for each fixed sufficiently large \(S\).  Indeed the unit tube over
this set lies in \(B_{2S+1}(Z)\cap B_{r+1}(y)\).

To prove \eqref{eq:B3}, first observe the exact identity
\[
 e(v)=\frac{q^2}{2}-\frac{q^2D}{4}.
\]
Together with Lemma~\ref{lem:finite-critical} this gives, uniformly at
bad centers,
\begin{equation}\label{eq:finite-qmass}
 R^{-m}\int_{B_R(y)}q^2\longrightarrow \sigma K\omega_m.
\end{equation}
Its radial distributional consequence is
\[
 R^{-m}\int q^2F(|x-y|/R)\,\dd x
 \longrightarrow
 \sigma K m\omega_m\int t^{m-1}F(t)\,\dd t.
\]
For completeness, the part of this integral at distance greater than
\(P\) from the zero set is \(o_P(1)R^m\), uniformly for \(R\ge1\).
Take a maximal unit-separated net on \(\Sigma\).  A fixed small ball
about any zero has energy bounded below: the uniform gradient bound
keeps \(|v|\le1/2\) in a uniform neighborhood of that zero.
Disjoint such balls and the density bound give at most
\(C(a+1)^m\) net points in any ball of radius \(a\).
The portion with zero-set distance in \([j,j+1]\), inside \(B_{2R}(y)\),
is therefore covered with volume at most
\[
 C(R+j+2)^m(j+2)^n.
\]
Exponential decay \eqref{eq:phase-decay}, summed over \(j\ge P-1\),
proves the assertion.

For fixed \(P\), points within \(P\) of a zero with \(s\le2S\)
have volume \(o(R^{m-2})\) in \(B_{2R}(y)\), by the fixed enlargement
of the bad set.  Every remaining point within \(P\) of \(\Sigma\)
has a nearest zero in \(\{\chi=1\}\), and belongs to the injective
good tube.  Replacing \(F(|\Phi(z,t)-y|/R)\) by \(F(|z-y|/R)\)
costs \(O(P/R)\) in normalized mass.  The normal profile integral
is at most \(\sigma+o_S(1)\).  First choose \(P\) to make the discarded
tail small, then \(S\) to make this profile and Jacobian error small.
To make the direction of this transfer explicit, the preceding estimates
give
\[
 R^{-m}\int q^2F(|x-y|/R)\,\dd x
 \le (\sigma+o_S(1))R^{-m}\int_\Sigma\chi F(r/R)\,\dd S
       +o_P(1)+o_{R\to\infty;P,S}(1).
\]
The error \(o_P(1)\) tends to zero as \(P\to\infty\); with \(P\)
fixed, \(o_S(1)\) tends to zero as \(S\to\infty\).
The radial mass limit therefore gives \eqref{eq:B3} in precisely the
stated order of choices.

\emph{The Hardy coefficient.}
The radial $q^2$ mass has now been identified, including its coefficient.
Put \(a=(m-2)/2\). In \eqref{eq:q-stability} use
\(\eta(x)=|x-y|^{-a}\) on \(R_0/2<|x-y|<4R\), extended with
smooth cutoffs on the adjacent inner and outer dyadic shells.
The cutoff-shell costs are \(O(1)\), uniformly in \(y\) and \(R\);
this follows from \(\int_{B_t(y)}q^2\le Ct^m\).
The remaining right-hand integral is
\[
 a^2\int_{\{R_0/2<|x-y|<4R\}}q^2|x-y|^{-m}\,\dd x.
\]
Stieltjes integration by parts in \eqref{eq:finite-qmass} gives
\[
 \int_{\{b<|x-y|<cR\}}q^2|x-y|^{-m}\,\dd x
 =\sigma K m\omega_m\log R+o(\log R)
\]
for fixed \(b,c>0\).  The same leading term holds after including
the cutoff shells.  Division by \(\sigma\log R\) proves
\eqref{eq:B4}.
\end{proof}

\subsection{Nearest-sheet distance and its weak Laplacian}

The mass estimates are now available.  To use them in a Green test, we
need an upper Laplacian bound for the distance to a neighboring sheet
and an exact count of how often a target point can receive a curvature
term.  The following argument works with local adjacent sheets and does
not require a global labeling of the zero set.

\begin{lemma}[Nearest feet and multiplicity]\label{lem:finite-gap}
On \(\Sigma_S\), for \(S\) sufficiently large, let \(l(z)\) be the
shortest Euclidean distance to a distinct sheet in the local good chart.
This definition is independent of the chart.  Uniformly as \(s(z)\to\infty\),
\[
 l(z)\longrightarrow\infty,\qquad l(z)=o(s(z)).
\]
Locally \(l\) is the minimum of the smooth distance functions to the
one or two adjacent sheets.  On an active branch \(l_b\), with nearest
foot \(z'=\pi_b(z)\), one has
\begin{equation}\label{eq:B5}
 \Delta_\Sigma l_b
 \le |H(z)|+|H(z')|+\frac{C}{s(z)}|\nabla_\Sigma l_b|^2.
\end{equation}
Every measurable choice of an active foot has inverse multiplicity at
most two.  The branch area Jacobian is \(1+o_S(1)\), and
\[
 s(z')/s(z)=1+o_S(1),\qquad
 |z'-y|/|z-y|=1+o_S(1)\quad(y\in Z).
\]
\end{lemma}

\begin{proof}
\emph{Defining the adjacent foot.}
All \(K\ge2\) graphs have height \(o(s)\); their separation diverges,
so the nearest distinct-sheet distance has the claimed size.
Orientations alternate because the solution changes sign on crossing
each regular zero graph.  Equivalently, \(l\) is the nearest distance
to a zero whose oriented normal has scalar product less than
\(-1/2\) with \(\nu(z)\).  The minimum is attained within \(o(s)\),
and every competitor that could improve it lies in the same buffered
chart.  A nonadjacent sheet cannot minimize: the chord to it crosses
an intervening zero sheet at a strictly smaller distance.
This also proves that the definition agrees on overlapping charts.

Projection to an adjacent sheet is smooth at these distances.
Indeed its curvature is \(O(s^{-1})\), the distance is \(o(s)\), and
the slopes are small.  In graph coordinates the squared-distance
function has positive definite Hessian near its minimizing foot;
all possible minimizers lie in this region, since their horizontal
displacements are at most \(l=o(s)\).  It thus has a unique foot.

\emph{The branch Laplacian.}
Choose on the target graph the unit normal \(N\) pointing toward the
source.  Then \(z=z'+l_bN(z')\), and the distance Hessian at \(z\)
has eigenvalues
\[
 \frac{\kappa_i}{1+l_b\kappa_i}\le\kappa_i
\]
in the transported target tangent directions, and zero in the normal
direction.  Here \(1+l_b\kappa_i>0\); its operator norm is \(O(s^{-1})\).
The inequality follows exactly by subtracting
\(l_b\kappa_i^2/(1+l_b\kappa_i)\).
Restricting to the source surface removes its normal-normal Hessian
component.  Its absolute value is bounded by
\[
 \frac Cs\bigl(1-(\nu(z)\cdot N)^2\bigr)
   =\frac Cs|\nabla_\Sigma l_b|^2.
\]
The restriction formula for the Laplacian adds
\(-H(z)\nu(z)\cdot N\), of size at most \(|H(z)|\).
Taking the ambient trace and bounding the target trace by
\(|H(z')|\) proves \eqref{eq:B5}.

\emph{The Jacobian and the count of two.}
The exact projection area Jacobian is
\begin{equation}\label{eq:finite-foot-jacobian}
 J_{\pi_b}(z)
  =\frac{|\nu(z)\cdot N(z')|}
         {|\det(\Id+l_b A_N(z'))|}
  =1+o_S(1).
\end{equation}
To check multiplicity globally, fix a target point \(z'\).
Write \(l(z)\le\varepsilon_S s(z)\) on \(\Sigma_S\), where
\(\varepsilon_S\to0\).  The Lipschitz property of \(s\) gives
\[
 |z-z'|\le\frac{\varepsilon_S}{1-\varepsilon_S}s(z').
\]
Choose \(S\) so this fraction is smaller than the fixed inner radius
of a good chart.  Every active source projecting to \(z'\) then
lies in that single chart centered at \(z'\).
The fixed normal line through \(z'\) meets each small-slope graph
at most once.  Any intervening sheet would meet the source-target
segment and contradict minimality.  Only the immediately adjacent
sheet on either side can therefore supply a source, giving at most
two preimages.  This reasoning also covers ties between active branches.
A countable atlas and a measurable selection among the finite local
branches permit the area formula with this multiplicity.
Finally \(s\) and distance to a fixed \(y\in Z\) are 1-Lipschitz and
\(s(z)\le|z-y|\); the two ratio estimates follow from \(l=o(s)\).
\end{proof}

Let \(\psi\) be nonnegative, increasing and concave on the range of \(l\), and set
\(w=\min\{T_0,\psi(l)\}\).  If
\[
 -\psi''(l)/\psi'(l)\ge C/s,\qquad 0\le\psi'(l)\le L,
\]
Lemma~\ref{lem:finite-gap} gives
\begin{equation}\label{eq:B6}
 \Delta_\Sigma w\le L\bigl(|H(z)|+|H(\pi(z))|\bigr)
 \quad\hbox{on }\Sigma_S
\end{equation}
in the sense of measures, where any measurable active foot may be
chosen almost everywhere.  To justify the weak statement, each local
distance branch is smooth with locally bounded second derivatives.
Their finite minimum, including the constant cap, is locally semiconcave.
Its singular Hessian part is nonpositive, and its absolutely continuous
Hessian agrees almost everywhere with that of an active branch.
On the capped part the constant branch has zero Laplacian and the
right-hand side is nonnegative.  The chain rule therefore gives
\eqref{eq:B6}.  In particular \(\psi(l)=\sqrt l\) works with \(L=1\)
for large \(S\), since \(l=o(s)\).

\subsection{A Green test and a first bound on the gap}

The concave transformation in \eqref{eq:B6} absorbs the gradient error
in the distance Laplacian.  The next test integrates this inequality
against an inward Green kernel.  Its cutoff errors must be independent
of the cap \(T_0\), except for a geometric tail whose coefficient can
be made small before the outer radius is increased.

Fix henceforth \(F\in C_c^\infty((1/2,2))\), nonnegative and positive
on a neighborhood of \(1\), normalized by
\(\int t^{m-1}F(t)\,\dd t=m-2\).
Define the radial function \(\phi_R\), for \(0<r<2R\), by
\[
 \phi_R(r)=\int_r^{2R}t^{1-m}
       \left(R^{-m}\int_t^{2R}u^{m-1}F(u/R)\,\dd u\right)\dd t,
\]
and set it to zero for \(r\ge2R\).  Direct differentiation gives
\[
 \phi_R''+\frac{m-1}{r}\phi_R'=R^{-m}F(r/R),\qquad
 0\le\phi_R\le r^{2-m},\qquad
 r|\phi_R'|+r^2|\phi_R''|\le Cr^{2-m}.
\]

\begin{lemma}[Green inequality with cap-independent collar errors]
\label{lem:finite-green}
Suppose that \eqref{eq:B6} holds, with \(L\le2\), and that
\(\psi(l)\) is bounded uniformly on each fixed range \(S<s<A\).
For fixed sufficiently large \(S\) and \(R_0\ge10\), all caps \(T_0>0\),
all \(y\in Z\), and large \(R\), one has
\begin{equation}\label{eq:B7}
 \begin{split}
 R^{-m}\int_\Sigma w\chi F(r/R)\,\dd S
 &\le L\int_{\Sigma_S\cap\{R_0<r<2R\}}
       r^{2-m}\bigl(|H(z)|+|H(\pi(z))|\bigr)\,\dd S\\
 &\quad+C_{R_0,S}+o_{R\to\infty;R_0,S}(\log R)
       +T_0\eta(R_0).
 \end{split}
\end{equation}
Here \(C_{R_0,S}\) is independent of \(T_0,y,R\). The
\(o(\log R)\) remainder is independent of \(T_0\) and uniform in \(y\);
\(\eta(R_0)\to0\) is uniform in \(y\).
These constants may depend on the fixed function \(\psi\).
\end{lemma}

\begin{proof}
The cap can affect the radial geometric error, but it must not multiply
the errors near the fixed inner boundary or the bad set. We keep those
three terms separate.
Choose a smooth radial \(\alpha\), zero on \(r\le R_0\), one on
\(r\ge2R_0\), with derivatives bounded by \(C R_0^{-j}\).
Test \eqref{eq:B6} against the nonnegative compactly supported
\(\alpha\chi\phi_R\).  Its support lies in the regular open surface,
so distributional integration by parts has no singular-set boundary.
For a radial function \(f\), the convention \(H=\operatorname{div}_\Sigma\nu\) gives
\[
 \Delta_\Sigma f
 =f''+(m-1)f'/r
 +(f'/r-f'')(e_r\cdot\nu)^2-Hf'(e_r\cdot\nu).
\]
Thus the error between \(\Delta_\Sigma\phi_R\) and its flat
\(m\)-dimensional radial Laplacian is at most
\[
 C\bigl((e_r\cdot\nu)^2r^{-m}+|H|r^{1-m}\bigr).
\]
It contributes at most \(T_0\eta(R_0)\), by \eqref{eq:B2}.

Terms differentiating \(\alpha\) are supported where \(r\le2R_0\),
and there \(s\le r\).  The uniform bound \(l=o(s)\), after fixing
its lower threshold \(S\), implies \(w\le\psi(l)\le C_{R_0,S}\)
on this set, independently of the cap.  Area and derivative bounds
there give a cap-independent constant error.
Terms differentiating \(\chi\) are supported in \(S<s<2S\), where
again \(w\le C_S\), independently of the cap.  The ambient derivative
bounds for \(\chi\) and the bounded mean curvature imply that these
terms are bounded in absolute value by
\[
 C_S\int_{\Sigma\cap\{S<s<2S,\ R_0<r<2R\}}
                   r^{2-m}\,\dd S.
\]
On a dyadic shell of radius \(a\), \eqref{eq:finite-collar-area}
bounds this by a quantity tending to zero as \(a\to\infty\),
uniformly in \(y\).  The sum of \(O(\log R)\) such quantities is
\(o(\log R)\): split at any fixed large shell, then use the uniform
small bound on all later shells.  This is the required Cesaro
estimate.

On the support of \(F(r/R)\), \(\alpha=1\) for large \(R\).
The main term on the left is therefore exactly that in
\eqref{eq:B7}.  Bound the nonnegative right-hand test by
\(r^{2-m}\); all remaining terms have just been estimated.
\end{proof}

\begin{lemma}[Rough gap bound]\label{lem:finite-rough}
There is \(C<\infty\) such that, at all sufficiently distant good zeros,
\begin{equation}\label{eq:B8}
 l(z)\le C\bigl((\log s(z))^2+s(z)^{1/10}\bigr).
\end{equation}
\end{lemma}

\begin{proof}
Use $w=\min\{T_0,\sqrt l\}$. The Green inequality will bound its
average by $O(\log R)$. A Harnack estimate for adjacent height
differences then turns this average bound into a bound at any chosen
good zero.

\emph{The average estimate.}
The curvature integral in \eqref{eq:B7} is at most \(C(1+\log R)\).
Indeed \(|H|\le CD\), and each dyadic shell contributes \(O(1)\)
by \eqref{eq:B1}.  For its foot term use
\eqref{eq:finite-foot-jacobian}, multiplicity two and the comparable
radial weights; the image is contained in
\(\Sigma_{S/2}\cap\{R_0/2<r<4R\}\).
The same shell estimate applies.

\emph{From one gap to a patch of comparable gaps.}
Given a good zero \(z\), choose a nearest \(y\in Z\), and set
\(R=s(z)=|z-y|\), \(T_0=\sqrt{l(z)}\).
We claim that either \(l(z)\le C'R^{1/10}\), or an area-\(cR^m\)
patch of the source sheet, within distance \(cR\) of \(z\), has
\(l\ge c l(z)\).  In the chart of
Lemma~\ref{lem:finite-critical}, an adjacent height difference
\(h=f_{i+1}-f_i>0\) satisfies
\[
 \operatorname{div}(a(x')\nabla h)=\mathcal M(f_{i+1})-\mathcal M(f_i),\qquad
 \mathcal M(f)=\operatorname{div}\frac{\nabla f}{\sqrt{1+|\nabla f|^2}},
\]
where
\[
 a(x')=\int_0^1
 D_\xi\!\left(\frac{\xi}{\sqrt{1+|\xi|^2}}\right)
 \bigl(\nabla f_i+t(\nabla f_{i+1}-\nabla f_i)\bigr)\,\dd t.
\]
The small slopes make this matrix uniformly elliptic.  The right-hand
side has absolute value at most \(CR^{-2+1/10}\).
Interior Harnack for a positive solution with bounded forcing,
after rescaling to a unit disk, therefore gives on a fixed smaller disk
\[
 \inf h\ge c h(x'_z)-C R^{1/10}.
\]
Every adjacent vertical gap relevant to the source sheet is at least
\(l(z)\) at \(x'_z\).  If \(C'\) is sufficiently large, the asserted
lower bound follows for both adjacent gaps, or for the sole adjacent
gap on an extreme sheet.  Vertical gaps and nearest Euclidean distances
are comparable uniformly by the small slope bound; potential closest
points outside the buffered chart are at distance comparable to \(R\),
whereas \(l=o(R)\).  This proves the claim.

Shrink the patch fraction so \(F(r/R)\ge c>0\) on it.
On this patch \(s\ge(1-c)R\), hence \(\chi=1\) for large \(R\), and
\(w\ge cT_0\).  Thus the left side of \eqref{eq:B7} is at least \(c'T_0\).
Fix \(R_0\) large enough that \(\eta(R_0)<c'/2\), independently of
\(z,T_0,R\).  Absorb its last term and obtain
\(T_0\le C(1+\log R)\).  Squaring proves \eqref{eq:B8} in the second
case, and its first case already has the asserted bound.
\end{proof}

\subsection{Completion of the finite-density proof}

The rough estimate has strengthened \(l=o(s)\) to \(l^2=o(s)\).
This permits a concave transformation whose derivative tends to one,
so the coefficient of the mean-curvature term survives the Green test.
We first obtain the upper logarithmic coefficient and then prove that
almost all the surface mass attains the logarithmic cap.

\begin{proof}[Completion of the proof of Theorem~\ref{thm:finite}]
\emph{The nearly linear gap test.}
Fix a small tolerance \(\delta>0\) and use
\[
 \psi(l)=l+\log l,\qquad w=\min\{T_0,\psi(l)\}.
\]
Because \(l\to\infty\), \(\psi'\le1+\delta\) once \(S\) is large.
Moreover
\[
 -\frac{\psi''(l)}{\psi'(l)}=\frac1{l(l+1)},\qquad
 \frac{l(l+1)}{s}\longrightarrow0
\]
by \eqref{eq:B8}.  Thus \eqref{eq:B6} and Lemma~\ref{lem:finite-green}
apply with \(L\le1+\delta\).

\emph{The upper logarithmic coefficient.}
Let \(1/(2\sqrt2)<c_*<\sqrt2/3\) and \(P\) be those of
Lemma~\ref{lem:diffuse-mean}.  Increase \(S\) so this lemma applies on
\(\Sigma_{S/2}\) and its \(P\)-tube is injective.  The source term of
\eqref{eq:B7} contributes one curvature integral; the target terms
contribute at most two by Lemma~\ref{lem:finite-gap}.  The area Jacobians
and radial-weight changes cost a relative \(1+C\delta\) after increasing
\(S\).  Thus the curvature integral in \eqref{eq:B7}, including its
factor \(L\), is at most
\[
 (1+C\delta)\,3\int_{\Sigma_{S/2}\cap\{3R_0/4<r<3R\}}
                        |H(z)|r^{2-m}\,\dd S.
\]
At every such base point,
\[
 |H(z)|\le\frac{c_*}{\sigma}
             \int_{-P}^{P}q^3D(z+t\nu(z))\,\dd t.
\]
Choose \(R_0\) also large compared with \(P\).
Tube changes of radial weights and volume Jacobians have relative
error as small as desired, and this enlarged region's \(P\)-tube lies in
\(R_0/2<|x-y|<4R\) when \(R_0>4P\).
Using \eqref{eq:B4} therefore gives
\begin{equation}\label{eq:B9}
 \limsup_{R\to\infty}\frac{L}{\log R}
 \int_{\Sigma_S\cap\{R_0<r<2R\}}
 r^{2-m}\bigl(|H(z)|+|H(\pi(z))|\bigr)\,\dd S
 \le (1+C\delta)\,3c_*K m\omega_m\frac{(m-2)^2}{4}.
\end{equation}
There is no extra multiplicity from this last transfer: the normal
tube over \(\Sigma_{S/2}\) is injective.

\emph{Almost all surface mass reaches the cap.}
For the lower bound choose
\(T_0=(\sqrt2-\delta)\log R\).
We first prove that, for any \(\gamma>0\),
\begin{equation}\label{eq:finite-endpoint-defect}
 D(z)+D(\pi(z))\ge c_\gamma
                    e^{-(\sqrt2+\gamma)l(z)}
 \quad(z\in\Sigma_S)
\end{equation}
when \(S\) is large.
The endpoint normals have opposite orientations and \(D=o(1)\), so
\(|p(z)-p(\pi(z))|\ge c>0\).
Apply \eqref{eq:diffuse-local} with exponential rate
\(2\sqrt2+2\gamma\) along each half of the chord.  Since
\(|\nabla p|=|T|\le C\sqrt D\), integration from the two endpoints
to the midpoint gives
\[
 c\le C_\gamma
       \bigl(\sqrt{D(z)}+\sqrt{D(\pi(z))}\bigr)
        e^{(\sqrt2+\gamma)l(z)/2}.
\]
Squaring proves \eqref{eq:finite-endpoint-defect}.

Choose \(\gamma>0\) so
\(\beta=(\sqrt2+\gamma)(\sqrt2-\delta)<2\).
On the exceptional set
\[
 E_R=\{z\in\Sigma_S:R/2<r<2R,\ l(z)<T_0\},
\]
\eqref{eq:finite-endpoint-defect} gives
\(D(z)+D(\pi(z))\ge cR^{-\beta}\).
By \eqref{eq:B1}, the area formula, multiplicity two and the radial
comparability of the feet,
\[
 \cH^m(E_R)
 \le CR^\beta\int_{E_R}\bigl(D(z)+D(\pi(z))\bigr)\,\dd S
 \le CR^{m-2+\beta}=o(R^m).
\]
Outside this exceptional set on the support of \(F\), \(l\ge T_0\);
since \(l>1\), \(w=T_0\).  Applying \eqref{eq:B3}, with its error
chosen at most \(\delta\), gives
\[
 \liminf_{R\to\infty}\frac{R^{-m}}{\log R}
       \int_\Sigma w\chi F(r/R)\,\dd S
 \ge(\sqrt2-\delta)\bigl(Km\omega_m(m-2)-\delta\bigr).
\]
\emph{The incompatible coefficients.}
All cutoff errors in \eqref{eq:B7} other than \(T_0\eta(R_0)\)
vanish after division by \(\log R\).
To make the order of choices explicit, first fix \(\delta\);
then fix \(P,S\) to meet the mean-curvature, tube, concavity and mass
requirements; next fix \(R_0\), as large as needed for the radial
tail; finally let \(R\to\infty\).
Letting \(R_0\to\infty\) in the resulting numerical inequality
removes \((\sqrt2-\delta)\eta(R_0)\).
Now send \(\delta\to0\) and use \eqref{eq:B9}.  Cancellation of the
positive factor \(K m\omega_m(m-2)\) yields
\[
 \sqrt2\le \frac{3c_*}{4}(m-2).
\]
This is impossible for \(3\le m\le6\), since its right side is at most
\(3c_*<\sqrt2\).  In the endpoint \(m=6\), the strict inequality in
Lemma~\ref{lem:diffuse-mean} is essential.
No critical solution exists, proving Theorem~\ref{thm:finite}.
\end{proof}

\subsection{The local regularity consequence}

The classification has now made every finite threshold subcritical.
The following consequence therefore applies a local density bound at
one large radius; it does not require a global density bound for the
particular entire solution to which it is applied.

\begin{corollary}[Uniform regularity from any finite density bound]
\label{cor:finite-subcritical}
Fix \(4\le n\le7\) and \(0<K_f<\infty\).
There are \(c>0,C<\infty,r_0<\infty\), depending only on \(n,K_f\),
such that every entire stable parameter-one solution with
\(M_r(z)\le K_f\), \(r\ge r_0\), satisfies
\[
 |\nabla v|\ge C^{-1},\qquad
 \left|\nabla\frac{\nabla v}{|\nabla v|}\right|\le C/r
 \quad\hbox{on }B_{cr}(z)\cap\{|v|\le0.9\}.
\]
For each prescribed slope tolerance and each fixed \(\iota>0\),
if \(v(z)=0\), a sufficiently thin slab confinement of the zero set in
this ball about a hyperplane through \(z\) gives, on a smaller cylinder, ordered graphs with that
slope tolerance and
\[
 \|\nabla^2 f_i\|_\infty+
 r^{1/2}[\nabla^2 f_i]_{C^{1/2}}\le C'/r,\qquad
 |H_i|\le C'_\iota r^{-2+\iota}.
\]
The constants and interior fractions depend only on the stated
parameters.  On each fixed window about an inner zero, the solution
converges smoothly to its oriented planar profile as \(r\to\infty\),
uniformly among these solutions.
\end{corollary}

\begin{proof}
Theorem~\ref{thm:finite} makes every finite \(K_f\) subcritical.
Apply \cite[Theorem~3.2]{FSS2025} with its tolerance \(\delta(K_f,n)\)
and \(r^{-1}\le\delta\); its allowed density is \(K_f+\delta\),
so the stated bound suffices.  Its sheeting assumptions and
\cite[Theorem~2.14, Theorem~4.8, and Lemmas~4.9--4.10]{FSS2025}
give the remaining assertions, as specified above.  Constants
\(\pm1\) have an empty core band and cause no exception.
\end{proof}
